\chapter{Marco Teórico}

En esta sección se presentan los conceptos fundamentales relacionados con el reconocimiento de imagenes enfocado al 
diagnóstico de enfermedades foliares utilizando arquitecturas ligeras de aprendizaje profundo en dispositivos móviles.
Se abordan temas  los modelos de aprendizaje profundo, transfer learning y los desafíos asociados a la implementación
en dispositivos móviles.



\section{Redes neuronales convolucionales para clasificación de imágenes}
Las redes neuronales convulucionales (CNN) son un tipo de arquitectura de aprendizaje profundo que se ha vuelto muy popular en los 
ultimos años, especialemente en tareas de clasificación de imágenes. Estan diseñadas para reconocer objetos usando capas
que incluyen la convolución, la activación y el pooling. Estas capas permiten a las CNN aprender características jerárquicas de
las imágenes, desde bordes y texturas simples en las primeras capas hasta patrones más complejos en las capas posteriores.
Ademas estas capas trabajan juntas usando backprogation para adaptar y mejorar la red. El objetivo principal de esta las CNN es 
crear una red más extensa con menos parámetros \cite{agronomy13040961}.


\subsection{Transfer Learning y ajuste fino}
Las redes neuronales convolucionales (CNN) han demostrado un rendimiento sobresaliente en tareas de 
clasificación de imágenes, pero su entrenamiento desde cero requiere grandes cantidades de datos y recursos computacionales.
El transfer learning es una técnica que permite aprovechar el conocimiento adquirido por un modelo preentrenado en una tarea relacionada 
para mejorar el rendimiento en una nueva tarea con menos datos. En este enfoque, se utiliza un modelo preentrenado 
como punto de partida y se ajusta (fine-tuning) para adaptarlo a la nueva tarea específica, lo que puede resultar en una mejora
significativa en la precisión del modelo con un menor esfuerzo de entrenamiento.

La tecnica de transfer learning se ha vuelto especialmente relevante en el contexto de dispositivos móviles, donde las limitaciones
de recursos hacen que el entrenamiento desde cero sea impráctico. Hay una manera de acelerar el tiempo de entrenamiento, previniendo el 
overfitting, mejorando la interpretabilidad del modelo, reduciendo los requermientos de memoria y optimizando la eficiencia es utilizar
la tecnica de congelado de capas (layer freezing). Esta técnica consiste en mantener fijas las capas iniciales del modelo preentrenado,
ue suelen aprender características generales de las imágenes, y solo ajustar las capas finales, que se encargan de la
clasificación específica de la nueva tarea. Esto permite aprovechar el conocimiento general del modelo preentrenado mientras
se adapta a las particularidades de la nueva tarea con un menor riesgo de sobreajuste y una reducción significativa en el 
tiempo de entrenamiento \cite{zainab2025plant}.




\section{Arquitecturas ligeras para dispositivos móviles}

\section{Metodos de despligue}
Para hablar de la inferencia de modelos de inteligencias artificial en dispositivos móviles, es importante entender los conceptos de
inferencia en la nube e inferencia local. El paradigma tradicional de inferencia en la nube consisten en recolactar la informacion requerida
para la inferencia a nivel dispositivo El dispositivo se conecta a un servidor en la nube a traves de la red. Luego los modelos de inteligencia 
artificial operando en los servidores ejecutan la inferencia sobre la data y envian el resultado al dispositivo \cite{li2025empirical}.  

Los sitemas basados en la nube tienen acceso casi que a recursos ilimitados, unicamente limtados por el costo de los recursos utilizados, lo que los
hace ideales para escenarios de complejos de actividades de alto consumo de recursos como almacenamiento de data, procesamiento de data, 
entrenamiento de modelos, etc. Sin embargo, la inferencia en la nube tiene limitaciones que se vuelven mas promientes a medida IA constanteme se 
esta transformando en un factor clave en la vida cotidiana, especiamente den las aplicaciones móviles.\cite{dhar2019ondevice}. A continuacion se presentan algunas de las limitaciones mas importantes de la inferencia en la nube:
algunas de las desvantajas de la inferencia en la nube incluyen:

\begin{description}
    \item[Privacidad y seguridad:] La segurdidad es una de las preocupaciones mas importantes, ya que la data del usuario debe ser enviada a
servidores externos para su procesamiento y almacenamiento, lo que puede exponer la información a riesgos de seguridad y privacidad.
    \item[Latencia:] Esta es la descripción detallada del segundo punto.
\end{description}

Por otro lado la inferencia local, también conocida como inferencia en el dispositivo (on-device), se refiere a la ejecución de modelos 
de inteligencia artificial directamente en el hardware del dispositivo móvil, sin necesidad de enviar datos a servidores externos \cite{dhar2019ondevice}. Este enfoque tiene varias ventajas, como la reducción de la latencia, 
la mejora de la privacidad y la capacidad de funcionar sin conexión a internet lo que garantiza respuestas inmediatas y una mejora significativa
en la privacidad y seguridad de las aplicaciones. Sin embargo, también presenta desafíos técnicos
significativos debido a las limitaciones de recursos de los dispositivos móviles, como la memoria, el procesamiento y el consumo energético.
    
\subsection{Dispositivos edge}
Para enteder lo que implica la inferecias en el dispositivo se vuelve imperativo de definir lo que es un dispositvo edge. este se puede definir como
definir como aquel cuyos recursos de computación, memoria y energía son limitados y no pueden aumentarse ni disminuirse fácilmente ya sea por factor 
de forma o costo. dentro de esta categoría se encuentran los dispositivos móviles, como teléfonos inteligentes y tabletas,
así como también dispositivos IoT (Internet de las cosas)y otros dispositivos embebidos.

El hardware de los dispositivos edge se encuentra como la base de sus limitaciones. Inclusive es una linea de investigacion con el objetivo de 
desarrollar hardware especializado para mejorar el rendimiento de la inferencia en el dispositivo, como los procesadores de inteligencia
artificial (AI accelerators) y las unidades de procesamiento gráfico (GPUs) optimizadas para tareas de aprendizaje profundo. Sin embargo, 
es un proceso es un proceso costoso que requiere una gran inversión de capital para construir laboratorios e instalaciones de fabricación, y suele implicar
largos plazos de tiempo para su desarrollo.



\subsection{TensorFlow Lite, LiteRT y cuantización}
Todos los algoritmos de aprendizaje automático dependen de operaciones clave como la multiplicación-acumulación o Multiply-Add
en el caso de las redes neuronales. Las bibliotecas que soportan estas operaciones constituyen la interfaz que separa el hardware de 
los algoritmos de aprendizaje. Esta separación permite el desarrollo de algoritmos que no dependen de ninguna arquitectura de hardware específica.

Esta separación entre el software y el hardware es fundamental para la flexibilidad en la inteligencia artificial,
ya que permite que un mismo modelo pueda ejecutarse en diferentes dispositivos sin necesidad de ser rediseñado desde cero.

\section{Desafios tecnicos en la implementación de modelos de aprendizaje profundo en dispositivos móviles}

\subsection{Tiempo de procesamiento}
El tiempo de respuesta suele estar entre los factores más críticos para la usabilidad de cualquier aplicación en el dispositivo (on-device). 
Las dos medidas comúnmente utilizadas son el rendimiento (throughput) y la latencia. El rendimiento se mide como la tasa a la que se procesan los datos de 
entrada. Para maximizar el rendimiento, es común agrupar las entradas en lotes (batches), lo que resulta en una mayor utilización.
Sin embargo, realizar esta medición conlleva un tiempo de espera adicional para agregar los datos en dichos lotes. 
Por tanto, para casos de uso críticos en tiempo, la latencia es la medida utilizada con mayor frecuencia.

 \subsection{Memoria}
La memoria de un dispositivo proporciona un acceso inmediato a los datos en comparación con otros componentes de almacenamiento. 
Sin embargo, esta velocidad de acceso tiene costos más elevados. Por razones de costo, la mayoría de los dispositivos edge están diseñados con una memoria 
limitada. Como tal, la huella de memoria (memory footprint) de las aplicaciones edge suele estar adaptada para un dispositivo objetivo específico.

Los algoritmos de aprendizaje automático avanzados a menudo consumen una cantidad significativa de memoria durante la construcción del modelo,
debido al almacenamiento de los parámetros del modelo, variables auxiliares, etc. Por ejemplo, incluso modelos de clasificación de imágenes relativamente 
simples como ResNet-50 pueden ocupar megabytes de espacio en memoria

\subsection{Consumo energético}
El consumo de energía es otro factor crucial para la ejecución de modelos en dispositivos móviles. Una solución energéticamente eficiente puede prolongar 
la vida útil de la batería y reducir los costos de mantenimiento. La potencia del sistema, comúnmente estudiada en el desarrollo de hardware, 
es la relación entre el consumo de energía y el lapso de tiempo para una tarea determinada. 
Vincular el consumo de energía de una aplicación particular impulsada por IA para un dispositivo específico depende conjuntamente de una
serie de variables como el tiempo de ejecución, la memoria, etc. Capturar estas dependencias casi nunca es determinista.
Por ello, la mayoría de las investigaciones existentes estiman el uso de potencia/energía a través de una función sustituta (surrogate function)
que típicamente depende de la memoria y el tiempo de ejecución de una aplicación.


